\documentclass[11pt]{article}
\usepackage[T1]{fontenc}
\usepackage{lmodern}
\usepackage[margin=1in]{geometry}
\usepackage{amsmath,amssymb,amsthm,mathtools}
\usepackage[hidelinks]{hyperref}
\usepackage{enumitem}
\usepackage{graphicx}
\emergencystretch=1.5em

\newtheorem{theorem}{Theorem}
\newtheorem{lemma}[theorem]{Lemma}
\newtheorem{corollary}[theorem]{Corollary}
\newtheorem*{remark}{Remark}
\newcommand{\F}{\mathbb F}
\newcommand{\Lp}{L_p[0,1]}
\newcommand{\B}{\mathrm B}
\newcommand{\Emb}{\operatorname{Emb}}
\newcommand{\GL}{\operatorname{GL}}
\newcommand{\E}{\mathbb E}

\title{Finite-\(p\) Extrinsic Metrics as\\Moment-Constrained Coupling Problems}
\author{Run \texttt{fa\_banach\_001}, lane 13}
\date{July 21, 2026}

\begin{document}
\maketitle

\begin{abstract}
Barto\v{s}ov\'{a}, Lopez-Abad, Lupini, and Mbombo close their paper by
saying that they do not know an explicit description of their three
extrinsic metrics for finite~$p$, analogous to their $p=\infty$ appendix.
We give exact descriptions for every $1\leq p<\infty$.  A finite family of
maps from finite-dimensional spaces into $L_p[0,1]$ is encoded exactly by
the joint law of its coefficient-vector functions.  Consequently the
full-rank metric is a moment-constrained optimal coupling problem, the
Kadets metric is a moment-constrained Hausdorff problem, and the square
rank-$k$ metric is the infimum of the norm of an explicit finite-rank
integral-kernel operator.  The formulas contain only probability laws on
finite-dimensional vector spaces and scalar $p$-moments.  They also apply
to the source's finite-support metrics through its Theorem~3.11.
\end{abstract}

\section{Source question and precise answer}

The source is D.~Barto\v{s}ov\'{a}, J.~Lopez-Abad, M.~Lupini, and
B.~Mbombo, \emph{The Ramsey properties for Grassmannians over
$\mathbb R,\mathbb C$}, arXiv:1910.00311v1 \cite{source}.  It defines three
extrinsic metrics: $\partial_{X,E}$ for norms represented by full-rank
embeddings, $\gamma_{k,E}$ for isometry types of $k$-dimensional subspaces,
and $d_{k,E}$ for $\GL(\F^k)$-orbits attached to rank-$k$ operators
$E^*\to E$.  After giving concrete descriptions when $p=\infty$, the
authors write on PDF page~24:
\begin{quote}
``We do not know a similar explicit description of the extrinsic metrics
for $p$'s other than $\infty$.''
\end{quote}
See Figure~\ref{fig:question}.

\begin{figure}[ht]
\centering
\includegraphics[width=.97\textwidth]{figures/open_problem_crop.png}
\caption{The closing question in arXiv:1910.00311v1, PDF p.~24.}
\label{fig:question}
\end{figure}

We interpret ``description'' as an exact variational formula that eliminates
the ambient infinite-dimensional embeddings.  Under that precise and
natural interpretation, Theorem~\ref{thm:formulas} gives a full answer for
all finite $p$.  It is a probabilistic finite-dimensional formula, rather
than a claim that the metrics always reduce to elementary closed forms.

\section{Proof intuition}

Choose a basis of a finite-dimensional space $W$.  Every bounded map
$A:W\to L_p[0,1]$ has the pointwise form
\[
       (Ax)(\omega)=a(\omega)(x),
       \qquad a(\omega)\in W^*.
\]
For several maps into the same $L_p$ space, retain the joint law of their
coefficient vectors.  The condition that $A$ is an isometric embedding is
then exactly a list of moment identities
\[
       \int |a(x)|^p\,d\mu(a)=\|x\|^p \quad(x\in W).
\]
Conversely, every such law can be realized by a random coefficient vector
on $[0,1]$.  Thus no information relevant to simultaneous embeddings is
lost when the maps are replaced by a probability coupling.

For $\partial$, the norm of the difference of two embeddings becomes the
supremum of one scalar $p$-moment.  For $\gamma$, the two subspace unit
balls become two coordinate unit balls inside the common seminorm induced
by the coupling, so their opening is an ordinary Hausdorff formula.  For
$d_k$, write $T=T_0T_1^*$.  If $a_0(\omega)\in V^*$ and
$a_1(\eta)\in V$ are the two coefficient vectors, then $T$ is represented
by the kernel $a_0(\omega)(a_1(\eta))$.  Subtracting two factorizations gives
the kernel in Theorem~\ref{thm:formulas}(c).

\section{Coefficient laws}

Throughout, $\F=\mathbb R$ or $\mathbb C$, $1\leq p<\infty$, and
$q\in(1,\infty]$ is conjugate to $p$.  All probability laws below are
Borel laws with the displayed finite $p$-moments.

\begin{lemma}[Joint coefficient-law lemma]\label{lem:law}
Let $W_1,\ldots,W_s$ be finite-dimensional $\F$-spaces.  Simultaneous
bounded maps $A_j:W_j\to L_p[0,1]$ determine a Borel probability law $\mu$
on $W_1^*\times\cdots\times W_s^*$ such that, if $z_j$ is the $j$th
coordinate,
\begin{equation}\label{eq:moment-map}
       \|A_jx\|_p^p=\int |z_j(x)|^p\,d\mu(z)
       \qquad(x\in W_j).
\end{equation}
Conversely, every such law is realized by simultaneous maps on
$L_p[0,1]$, with \eqref{eq:moment-map}.  Moreover, all joint pointwise
linear combinations of the maps have exactly the corresponding moments
under $\mu$.
\end{lemma}

\begin{proof}
For a basis $(e_i)$ of $W_j$, choose measurable representatives of
$A_je_i$ and put
\[
       a_j(\omega)=\sum_i (A_je_i)(\omega)e_i^*.
\]
Then $(A_jx)(\omega)=a_j(\omega)(x)$ outside one common null set.  The
joint push-forward of Lebesgue measure by $(a_1,\ldots,a_s)$ is $\mu$.
Finite dimensionality and the fact that every coordinate belongs to $L_p$
give the required finite moments, and integration proves
\eqref{eq:moment-map} and the final assertion.

Conversely, a Borel probability law on a finite-dimensional real or complex
vector space is the law of a Borel random variable on $[0,1]$ (the standard
randomization theorem; see, for example, \cite{Kallenberg}).  Realize the
coordinate tuple $z(\omega)$ and define
$(A_jx)(\omega)=z_j(\omega)(x)$.  The moment assumption puts these functions
in $L_p$, and \eqref{eq:moment-map} follows directly.
\end{proof}

Let $V$ be $k$-dimensional.  For norms $m,n$ on $V$, let
$\mathcal C_p(m,n)$ be the laws $\mu$ of pairs $(a,b)\in V^*\times V^*$
satisfying
\begin{equation}\label{eq:pair-constraints}
 \int |a(x)|^p\,d\mu=m(x)^p,
 \qquad
 \int |b(x)|^p\,d\mu=n(x)^p
 \quad(x\in V).
\end{equation}
Set
\[
 \rho_\mu(x,y)=
 \left(\int |a(x)-b(y)|^p\,d\mu(a,b)\right)^{1/p}.
\]
For a reference norm $r$ on $V$, define
\[
 D_r(\mu)=\sup_{r(x)\leq1}\rho_\mu(x,x),
\]
and define
\begin{align}
 G_{m,n}(\mu)=\max\bigg\{&
 \sup_{m(x)\leq1}\inf_{n(y)\leq1}\rho_\mu(x,y),\nonumber\\
 &\sup_{n(y)\leq1}\inf_{m(x)\leq1}\rho_\mu(x,y)
 \bigg\}.\label{eq:G}
\end{align}

For the square metric, a law $\mu$ is now on
\[
 \mathcal Z_V=V^*\times V\times V^*\times V,
 \qquad z=(a_0,a_1,b_0,b_1).
\]
It induces norms
\begin{align*}
 m_{0,\mu}(x)&=\left(\int|a_0(x)|^p\,d\mu\right)^{1/p},&
 m_{1,\mu}(f)&=\left(\int|f(a_1)|^p\,d\mu\right)^{1/p},\\
 n_{0,\mu}(x)&=\left(\int|b_0(x)|^p\,d\mu\right)^{1/p},&
 n_{1,\mu}(f)&=\left(\int|f(b_1)|^p\,d\mu\right)^{1/p}.
\end{align*}
For $\mathbf m,\mathbf n\in\mathcal D_k(L_p)$, let
$\mathcal Q_p(\mathbf m,\mathbf n)$ consist of those $\mu$ for which these
are genuine norms and
\[
 [(m_{0,\mu},m_{1,\mu})]=\mathbf m,
 \qquad [(n_{0,\mu},n_{1,\mu})]=\mathbf n
\]
in the source's diagonal $\GL(V)$-orbit space.  On two independent copies
$z=(a_0,a_1,b_0,b_1)$ and $z'=(a'_0,a'_1,b'_0,b'_1)$, put
\begin{equation}\label{eq:kernel}
 K_\mu(z,z')=a_0(a'_1)-b_0(b'_1)
\end{equation}
and
\begin{equation}\label{eq:kernelop}
 (\mathsf K_\mu g)(z)=\int K_\mu(z,z')g(z')\,d\mu(z').
\end{equation}
We use the complex bilinear realization of the duality
$L_p^*=L_q$, namely $\langle f,g\rangle=\int fg$.  With the usual
sesquilinear convention, one conjugates $g$ in \eqref{eq:kernelop}; the
operator norm is the same.

\begin{theorem}[Exact finite-$p$ descriptions]\label{thm:formulas}
For every $1\leq p<\infty$ the following hold.
\begin{enumerate}[label=\textup{(\alph*)}]
\item If $X=(V,r)$ and $m,n\in\mathcal N_V(L_p)$, then
\begin{equation}\label{eq:partial}
 \boxed{\displaystyle
 \partial_{X,L_p}(m,n)=\inf_{\mu\in\mathcal C_p(m,n)}D_r(\mu).}
\end{equation}
\item If $[m],[n]\in\mathcal B_k(L_p)$, then
\begin{equation}\label{eq:gamma}
 \boxed{\displaystyle
 \gamma_{k,L_p}([m],[n])=
 \inf_{\mu\in\mathcal C_p(m,n)}G_{m,n}(\mu).}
\end{equation}
\item If $\mathbf m,\mathbf n\in\mathcal D_k(L_p)$, then
\begin{equation}\label{eq:square}
 \boxed{\displaystyle
 d_{k,L_p}(\mathbf m,\mathbf n)=
 \inf_{\mu\in\mathcal Q_p(\mathbf m,\mathbf n)}
 \|\mathsf K_\mu\|_{L_q(\mu)\to L_p(\mu)}.}
\end{equation}
\end{enumerate}
The formulas are invariant under changing representatives in the indicated
orbit classes.
\end{theorem}

\begin{proof}
For (a), take simultaneous isometric embeddings
$T:(V,m)\to L_p$ and $U:(V,n)\to L_p$.  Lemma~\ref{lem:law} gives a law
$\mu\in\mathcal C_p(m,n)$, and pointwise
\[
       ((T-U)x)(\omega)=a(\omega)(x)-b(\omega)(x).
\]
Therefore $\|T-U\|_{(V,r)\to L_p}=D_r(\mu)$.  Conversely every admissible
$\mu$ produces exactly such simultaneous isometric embeddings.  Taking
the infimum on the two equivalent parameter spaces proves
\eqref{eq:partial}.

For (b), the unit balls in the ranges $T(V)$ and $U(V)$ are respectively
\[
       \{Tx:m(x)\leq1\},\qquad \{Uy:n(y)\leq1\}.
\]
Their ambient $L_p$ distance satisfies $\|Tx-Uy\|_p=\rho_\mu(x,y)$.
Thus their Hausdorff distance is exactly \eqref{eq:G}.  Lemma~\ref{lem:law}
again identifies all simultaneous realizations with all laws in
$\mathcal C_p(m,n)$, giving \eqref{eq:gamma}.  Replacing $m$ or $n$ by an
isometric representative merely reparametrizes $x$ or $y$, which also
proves the asserted invariance here.

For (c), first take admissible factorizations
\[
 T=T_0T_1^*,\qquad U=U_0U_1^*,
\]
where $T_0,U_0:V\to L_p$ and $T_1,U_1:V^*\to L_p$ are injective and induce
the prescribed orbit pairs.  Apply Lemma~\ref{lem:law} to the four maps.
Their joint law lies in $\mathcal Q_p(\mathbf m,\mathbf n)$.  In the
coefficient notation above, for $g\in L_q(\mu)$,
\[
 T_1^*g=\int a_1(z')g(z')\,d\mu(z'),
 \qquad
 (T_0T_1^*g)(z)=
 \int a_0(a'_1)g(z')\,d\mu(z').
\]
The analogous identity for $U$ shows that $T-U=\mathsf K_\mu$, including
equality of the $L_q\to L_p$ operator norms.

Conversely, any $\mu\in\mathcal Q_p(\mathbf m,\mathbf n)$ can be realized
on $[0,1]$.  Its four coordinates define $T_0,T_1,U_0,U_1$ by evaluation.
The moment norms being genuine norms makes all four maps injective.  In
finite dimension, the adjoints $T_1^*,U_1^*$ are onto $V$, so both products
have rank $k$, and the orbit constraints make the factorizations admissible
for the source definition.  Formula \eqref{eq:kernelop} again gives
$T-U=\mathsf K_\mu$.  Taking infima proves \eqref{eq:square}.  Finally,
the use of orbit constraints rather than fixed representatives makes the
$\GL(V)$ invariance in (c) automatic.
\end{proof}

\begin{corollary}[The source metrics]\label{cor:source}
Let $E_p=(\F^\infty,\|\cdot\|_p)$ and let $\widehat E_p=L_p[0,1]$ be its
Fra\"iss\'{e} limit in the notation of the source.  On the domains
$\mathcal N_V(E_p)$, $\mathcal B_k(E_p)$, and $\mathcal D_k(E_p)$, the
three $E_p$-extrinsic metrics have the respective descriptions
\eqref{eq:partial}, \eqref{eq:gamma}, and \eqref{eq:square}.
The same formulas describe the metrics on the full limit-space domains.
\end{corollary}

\begin{proof}
The source identifies the finite-$p$ limit space with $L_p[0,1]$ and proves
in its Theorem~3.11 that
\[
 \partial_{X,E_p}=\partial_{X,\widehat E_p},\qquad
 \gamma_{k,E_p}=\gamma_{k,\widehat E_p},\qquad
 d_{k,E_p}=d_{k,\widehat E_p}
\]
on the corresponding finite-support domains \cite[Theorem~3.11]{source}.
Apply Theorem~\ref{thm:formulas} with $\widehat E_p=L_p[0,1]$.
\end{proof}

\section{Verification and edge cases}

The proof is exact and does not depend on numerical evidence.  The script
\path{code/finite_atomic_checks.py} generates 100 random finite atomic
coefficient laws.  It independently constructs the corresponding weighted
$\ell_2$ embeddings and the square-operator matrices, then checks the moment
identities and the equality between the factorized operator and the kernel
operator.  The command
\begin{quote}\small\ttfamily
conda run --no-capture-output -n sandbox python \textbackslash\\
\hspace*{2em}code/finite\_atomic\_checks.py --trials 100
\end{quote}
completed with maximum residual below $10^{-12}$.  This is a regression
check of the formulas, not a proof.

For $p=1$, $q=\infty$ and the same finite-rank kernel calculation is valid;
no reflexivity is used.  For complex scalars, the bilinear and customary
sesquilinear $L_p$ duality identifications differ only by conjugation.  For
$1<p<\infty$ there are no additional restrictions.  Although the source's
Ramsey factorization theorems exclude certain even $p\geq4$, the metric
description itself uses only the coefficient-law lemma and therefore holds
for every finite $p$.

\section{Novelty, limitations, and review recommendation}

The bounded novelty check on July~21, 2026 searched the run registry,
solution and attempt indexes, the local parsed arXiv corpus, the exact source
sentence, the paper title together with ``extrinsic metrics'' and $L_p$, the
four authors' names with those terms, and ``$E$-Kadets metric'' with $L_p$.
The arXiv/web searches returned the original arXiv:1910.00311 and unrelated
or adjacent papers, but no later paper claiming these formulas or an answer
to the closing question.  Novelty confidence is therefore moderate, not
definitive.

The main limitation is terminological: the source does not formally define
what degree of simplification counts as an ``explicit description.''  The
present formulas are exact and eliminate optimization over maps into an
infinite-dimensional ambient space, replacing it by moment-constrained laws
on fixed finite-dimensional coefficient spaces.  They are not generally
closed elementary expressions in the input norms.

Status: \texttt{candidate full solution -- likely valid}, subject to human
review.  The most important review points are (i) whether this precise
variational meaning of ``explicit'' matches the authors' intended question,
(ii) the complex duality convention in the kernel formula, and (iii) the
application of the source's Theorem~3.11 to the three finite-support domains.

\begin{thebibliography}{9}
\bibitem{source}
D. Barto\v{s}ov\'{a}, J. Lopez-Abad, M. Lupini, and B. Mbombo,
\emph{The Ramsey properties for Grassmannians over $\mathbb R$, $\mathbb C$},
arXiv:1910.00311v1 (2019).

\bibitem{Kallenberg}
O. Kallenberg,
\emph{Foundations of Modern Probability}, 2nd ed., Springer, 2002.
\end{thebibliography}

\end{document}
